\subsection{Độ nhạy cảm với nhiễu Gaussian ở các biến đặc trưng}

\subsubsection{Hồi quy}

\tablename~\ref{tab:regression-sens-noise} trình bày các kết quả về độ nhạy (tỉ
lệ thay đổi tương đối khi chạy trên dữ liệu gốc) các chỉ số MSE, RMSE, MAE và
$R^2$ đối với các độ lệch chuẩn nhiễu khác nhau. Các chỉ số độ nhạy cảm càng
cao cho thấy kết quả càng tệ so với khi chạy trên dữ liệu gốc.

\import{\tablespath}{regression-sens-noise}

Khi tăng nhiễu lên $\sigma=0.1$, độ nhạy tương đối của tất cả các mô hình vẫn
còn khá thấp (đa số MSE/RMSE/MAE sens chỉ khoảng 0.2–3\%), cho thấy mô hình
chưa bị ảnh hưởng đáng kể bởi nhiễu nhỏ. Tuy nhiên, một số mô hình phi tuyến
như Polynomial và RBF bắt đầu xuất hiện độ nhạy cao hơn nhẹ so với các mô hình
tuyến tính.

Khi tăng lên $\sigma=0.5$, độ nhạy tăng rõ rệt. Các mô hình tuyến tính như OLS,
MBGD và WLS có mức tăng sai số tương đối khoảng 10–20\%, cho thấy chúng bắt đầu
bị ảnh hưởng rõ rệt dưới nhiễu trung bình. Các mô hình regularization như
Ridge, Lasso và Elastic Net có xu hướng ổn định hơn tương đối, với mức tăng sai
số thấp hơn nhóm tuyến tính thuần túy. Trong nhóm phi tuyến, Polynomial và
Fourier thể hiện độ nhạy cao, trong khi RBF chỉ tăng ở mức trung bình và vẫn
giữ được độ ổn định tương đối tốt hơn so với hai mô hình còn lại.

Ở mức nhiễu cao nhất $\sigma=1.0$, sự khác biệt trở nên rõ rệt hơn. Các mô hình
tuyến tính truyền thống tăng sai số tương đối lên khoảng 20–40\%. Các mô hình
phi tuyến có xu hướng nhạy hơn, tuy nhiên mức độ không đồng đều: Polynomial
tăng mạnh, trong khi RBF chỉ ở mức trung bình (ví dụ MSE sens khoảng 0.29),
cho thấy khả năng chống chịu nhiễu tốt hơn so với nhiều mô hình phi tuyến khác.
Đặc biệt, Fourier cho thấy mức độ biến động rất lớn (ví dụ MSE sens vượt
gần 1), cho thấy mô hình này rất nhạy với nhiễu đầu vào trong cấu hình hiện
tại. Ngược lại, các mô hình có regularization (Ridge, Lasso, Elastic Net) vẫn
giữ được mức tăng tương đối thấp hơn so với phần lớn các mô hình phi tuyến, dù
vẫn suy giảm đáng kể khi nhiễu lớn.

Tổng thể, kết quả cho thấy độ nhạy tăng theo $\sigma$ một cách nhất quán trên
tất cả mô hình. Các mô hình tuyến tính và regularization có xu hướng ổn định
hơn trong nhiễu thấp đến trung bình. Các mô hình phi tuyến đạt hiệu năng tốt
trên dữ liệu sạch nhưng nhạy hơn với nhiễu, tuy nhiên RBF là một trường hợp
đáng chú ý khi vẫn duy trì được độ ổn định tương đối ngay cả khi nhiễu lớn.

\subsubsection{Phân lớp}

\tablename~\ref{tab:classification-sens-noise} cho thấy độ nhạy cảm tương đối
của các mô hình phân lớp khi thêm nhiễu Gaussian vào dữ liệu đầu vào. Nhìn
chung, các giá trị sens có xu hướng tăng (dương hơn) khi $\sigma$ tăng, cho
thấy hiệu năng suy giảm theo mức nhiễu.

\import{\tablespath}{classification-sens-noise}


Ở mức nhiễu thấp $\sigma = 0.1$, đa số các mô hình vẫn khá ổn định. Softmax và
Weighted Softmax có sens dương nhỏ (khoảng 1--7\%), cho thấy hiệu năng chỉ suy
giảm nhẹ. LDA có một số chỉ số mang giá trị âm, tức là hiệu năng thậm chí cải
thiện nhẹ trong điều kiện nhiễu nhỏ, phản ánh tính ổn định cao. Ngược lại, QDA
có giá trị sens âm rất lớn về độ lớn (ví dụ Accuracy sens xấp xỉ $-1.54$), cho
thấy sự thay đổi cực mạnh và không ổn định (dù một phần là cải thiện, nhưng
biến động lớn là dấu hiệu kém bền vững).

Khi tăng lên $\sigma = 0.5$, xu hướng suy giảm trở nên rõ rệt hơn. Softmax có
sens dương lớn ở các chỉ số macro (khoảng 24--25\%), cho thấy hiệu năng giảm
đáng kể. Weighted Softmax vẫn duy trì sens thấp hơn, thể hiện khả năng chống
nhiễu tốt hơn. LDA tiếp tục ổn định ở Accuracy và Precision (sens gần 0 hoặc âm
nhẹ), nhưng Recall và F1 bắt đầu có sens dương đáng kể, cho thấy mô hình bị ảnh
hưởng khi xét đến cân bằng giữa các lớp. QDA tiếp tục có biến động rất lớn (cả
âm và dương), phản ánh sự thiếu ổn định nghiêm trọng.

Ở mức nhiễu cao $\sigma = 1.0$, sự khác biệt càng rõ. Softmax có sens dương lớn
(xấp xỉ 27--29\%), cho thấy suy giảm mạnh. Weighted Softmax vẫn kiểm soát tốt
hơn với sens thấp hơn đáng kể (xấp xỉ 7--13\%). LDA duy trì Accuracy và
Precision ổn định (sens gần 0), nhưng Recall và F1 có sens dương lớn, cho thấy
hiệu năng tổng thể vẫn bị ảnh hưởng. QDA tiếp tục là mô hình kém ổn định nhất
với các giá trị sens có độ lớn rất cao, cho thấy phản ứng mạnh và khó kiểm soát
trước nhiễu.

Tổng thể, Weighted Softmax là mô hình ổn định nhất (sens nhỏ và nhất quán), LDA
ổn định ở một số chỉ số nhưng không đồng đều, Softmax suy giảm rõ khi nhiễu
tăng, còn QDA có độ nhạy rất lớn và thiếu ổn định trong mọi mức nhiễu.
